\section{终点：数据会不会不重要}

访谈最后问：数据问题会不会有一天彻底不重要？谢晨的回答是，他越来越认为智能越强，对知识和数据越饥渴。人越优秀越想学习，AI 也可能如此。终局可能不是“数据消失”，而是从向外部数据学习，转向在仿真和自我环境中持续学习。

\subsection{评测是卡口}

对机器人来说，谢晨认为最关键的问题可能是评测规模化。没有规模化评测，就不知道模型是否真的变聪明。对大语言模型和 Agent 来说，也需要更高阶的评价指标，因为模型越强，旧 benchmark 越容易失效。

\begin{importantbox}{评测也是数据问题}
评测集不是外部裁判，而是训练和研究的反馈源。没有好的评测，数据闭环就无法判断方向。
\end{importantbox}

\subsection{自我学习与仿真终局}

如果 AI 能在仿真环境中设定目标、尝试、失败、修正、再尝试，数据工厂不会消失，而是形态变化：从采集人类数据，转向构造环境、目标和反馈，让 AI 在其中持续学习。

\begin{knowledgebox}{从 data factory 到 environment factory}
当外部数据边际收益下降，真正重要的可能是构造可学习环境。仿真、评测、reward 和失败反馈会成为新的数据工厂。
\end{knowledgebox}

\subsection{本章小结}

数据不会消失，它会从静态语料转向动态环境。越强的智能，越需要更高质量的反馈和更复杂的学习环境。


